\documentclass[12pt]{article}

\input{../../_assets/preambulo.tex}


\begin{document}

    % 1. Foto de fondo
    % 2. Título
    % 3. Encabezado Izquierdo
    % 4. Color de fondo
    % 5. Coord x del titulo
    % 6. Coord y del titulo
    % 7. Fecha

    
    \input{../../_assets/portada}
    \portadaExamen{ffccA4.jpg}{Ecuaciones\\Diferenciales II\\Examen XI}{Ecuaciones Diferenciales II. Examen XI}{MidnightBlue}{-8}{28}{2026}{}

    \begin{description}
        \item[Asignatura] Ecuaciones Diferenciales II.
        \item[Curso Académico] 2024/2025.
        \item[Grado] Doble Grado en Ingeniería Informática y Matemáticas.
        \item[Grupo] Único.
        \item[Profesor] Rafael Ortega Ríos.
        \item[Descripción] Segunda prueba de clase.
        \item[Fecha] 26 de mayo de 2025.
        % \item[Duración] ---.
    \end{description}
    \newpage


    % ------------------------------------
    
    \begin{ejercicio}
        Se denota por $\theta(t, h)$ a la solución del problema de valores iniciales
        \[
        \ddot{\theta} + \sen \theta = 0, \quad \theta(0) = 6\pi, \quad \dot{\theta}(0) = 2h.
        \]
        Calcula $\del{\theta}{h}(3\pi, 0)$. Se justificará la existencia de esta derivada.
    \end{ejercicio}

    \begin{ejercicio}
        Para cada $\veps \in \bb{R}$ se denota por $\left]\alpha(\veps), \omega(\veps)\right[$ al intervalo maximal de la solución del problema
        \[
        \dot{x} = \veps x^3, \quad x(0) = 1.
        \]
        Determina el conjunto $\cc{G} = \left\{(t, \veps) \in \bb{R}^2 : \alpha(\veps) < t < \omega(\veps)\right\}$. Se incluirá una representación gráfica de $\cc{G}$.
    \end{ejercicio}

    \begin{ejercicio}
        Se supone que $F : \bb{R}^3 \to \bb{R}^3$ es continua y cumple $F^{-1}(0) = \{0\}$. La sucesión $\{x_n\}_{n \geq 1}$ con $x_n \in \bb{R}^3$ cumple:
        \begin{enumerate}[label=\roman*)]
            \item $\|x_n\| \leq 7$
            \item $\|F(x_n)\| \leq \frac{1}{n}$.
        \end{enumerate}
        ¿Se puede asegurar que $\{x_n\}$ es convergente?
    \end{ejercicio}

    \begin{ejercicio}
        Para cada $n = 1, 2, \dots$ se denota por $x_n(t)$ a la solución del problema de Cauchy
        \[
        \dot{x} = n \sen\left(\frac{x}{n}\right), \quad x(0) = 1.
        \]
        Calcula, si existe, $\lim\limits_{n \to \infty} x_n\left(\nicefrac{1}{2}\right)$.
    \end{ejercicio}

    \begin{ejercicio}
        Dada una función $F : \bb{R}^3 \to \bb{R}^4$ de clase $C^1$, demuestra que existe una función continua $A : \bb{R}^3 \times \bb{R}^3 \to \bb{R}^{4 \times 3}$ para la que se cumple
        \[
        F(x) - F(y) = A(x, y)(x - y), \quad \forall x, y \in \bb{R}^3.
        \]
        Además, $A(x, x)$ es la matriz Jacobiana de $F$ en $x$.
    \end{ejercicio}

    \newpage
    \setcounter{ejercicio}{0} % Reiniciar contador de ejercicios
    % \noindent
    % \textbf{Solución.}

\end{document}